\chapter*{Résumé}
\addcontentsline{toc}{chapter}{Résumé}

Dans le cadre de la transformation numérique du \textbf{Groupement Sanitaire Territorial Tanger--Tétouan--Al Hoceïma (GST-TTA)}, ce projet porte sur la conception et le développement d'une \textbf{plateforme intelligente d'analyse conversationnelle de données de santé}. Son objectif est de faciliter l'exploitation des données hospitalières en permettant aux professionnels de santé et aux décideurs de formuler leurs demandes en langage naturel et d'obtenir des analyses statistiques, des visualisations et des interprétations adaptées à leurs besoins, sans nécessiter une maîtrise approfondie des outils traditionnels d'analyse de données.

La solution proposée repose sur une architecture modulaire combinant une interface conversationnelle, un orchestrateur agentique basé sur \textbf{LangGraph}, un moteur d'analyse déterministe et une infrastructure de modèles de langage locaux. Cette architecture établit une séparation entre l'intelligence chargée de comprendre et de planifier les demandes et les composants responsables de l'exécution effective des traitements. Le moteur d'analyse s'appuie notamment sur \textbf{pandas} et \textbf{DuckDB} afin de produire des résultats calculés de manière reproductible. Un mécanisme d'\textit{evidence} permet en outre de relier les réponses générées aux résultats effectivement obtenus lors de l'analyse, renforçant ainsi leur traçabilité et leur vérifiabilité.

La plateforme a été conçue avec une attention particulière portée à la \textbf{confidentialité et à la souveraineté des données}. L'utilisation de modèles de langage locaux permet de limiter l'exposition des données sensibles à des services externes, tandis que l'architecture modulaire facilite l'évolution et l'intégration de nouvelles capacités d'analyse.

La réalisation a suivi une démarche itérative fondée sur une approche Agile, depuis l'analyse des besoins jusqu'à la conception, l'implémentation et la validation de la plateforme. Les tests réalisés ont permis de vérifier le fonctionnement des principaux composants et leur adéquation avec les besoins identifiés. Certaines limites restent toutefois à traiter, notamment l'authentification et la gestion fine des droits d'accès, l'industrialisation des tests et l'extension des capacités analytiques. Ces éléments constituent des perspectives naturelles pour les évolutions futures de la solution.
\\
\\ 
\\ 
\textbf{Mots-clés :} Intelligence artificielle générative, analyse conversationnelle, données de santé, grands modèles de langage, agents intelligents, LangGraph, analyse de données, souveraineté des données.















\chapter*{Abstract}
\addcontentsline{toc}{chapter}{Abstract}

As part of the digital transformation of the \textbf{Tanger--Tétouan--Al Hoceïma Territorial Health Group (GST-TTA)}, this project focuses on the design and development of an \textbf{intelligent conversational health data analysis platform}. Its objective is to facilitate the exploitation of hospital data by enabling healthcare professionals and decision-makers to formulate requests in natural language and obtain statistical analyses, visualizations, and interpretations adapted to their needs, without requiring advanced expertise in traditional data analysis tools.

The proposed solution is based on a modular architecture combining a conversational interface, an agentic orchestrator built with \textbf{LangGraph}, a deterministic analysis engine, and a local language-model infrastructure. This architecture separates the intelligence responsible for understanding and planning user requests from the components responsible for executing the actual computations. The analysis engine relies in particular on \textbf{pandas} and \textbf{DuckDB} to produce reproducible results. An \textit{evidence} mechanism further links generated responses to the results effectively obtained during the analysis, thereby improving their traceability and verifiability.

Particular attention was given to \textbf{data confidentiality and sovereignty}. The use of local language models limits the exposure of sensitive data to external services, while the modular architecture facilitates future evolution and the integration of additional analytical capabilities.

The project was developed through an iterative Agile approach, covering requirements analysis, system design, implementation, and platform validation. The tests conducted verified the operation of the main components and their compliance with the identified requirements. Some limitations remain to be addressed, particularly authentication and fine-grained access control, test industrialization, and the extension of analytical capabilities. These aspects constitute natural directions for future development.
\\ 
\\ 
\\ 
\textbf{Keywords:} Generative artificial intelligence, conversational analytics, health data, large language models, intelligent agents, LangGraph, data analysis, data sovereignty.






% %%%%%%%%%%%%%%%%%%%%%%%%%%%%%%%%%%%%%%%%%%%%%%%%%%%%%%%%%%%%%%%%%%%%%%%%%%%%%%
% % الملخص
% %%%%%%%%%%%%%%%%%%%%%%%%%%%%%%%%%%%%%%%%%%%%%%%%%%%%%%%%%%%%%%%%%%%%%%%%%%%%%%

\chapter*{\textarabic{الملخص}}
\addcontentsline{toc}{chapter}{\textarabic{الملخص}}
\begin{Arabic}
في إطار التحول الرقمي الذي يشهده \textbf{التجمع الصحي الترابي طنجة--تطوان--الحسيمة (ATT-TSG)}، يهدف هذا المشروع إلى تصميم وتطوير \textbf{منصة ذكية للتحليل الحواري للبيانات الصحية}. تتمثل الغاية الأساسية في تسهيل استغلال البيانات الاستشفائية، من خلال تمكين مهنيي الصحة وصناع القرار من التفاعل مع مجموعات البيانات باستخدام اللغة الطبيعية، والحصول على تحليلات إحصائية موثوقة وتمثيلات بيانية ملائمة، مع ضمان سرية البيانات وسيادتها.

تعتمد المنصة على بنية معيارية تجمع بين واجهة مستخدم مبنية باستخدام \textbf{tiKetlevS}، وخلفية تعتمد على \textbf{IPAtsaF}، ومنسق للوكلاء الذكيين مبني على \textbf{hparGgnaL}، بالإضافة إلى محرك حتمي لتحليل البيانات يعتمد بشكل أساسي على \textbf{sadnap} و\textbf{BDkcuD}. وتفصل هذه البنية بين دور نموذج اللغة ودور المحرك التحليلي؛ إذ يتولى النموذج فهم طلب المستخدم وتنسيق مراحل المعالجة، بينما ينفذ محرك التحليل العمليات الحسابية بطريقة حتمية وقابلة للتحقق. كما يتيح نظام \textit{ecnedive} ربط الإجابات المقدمة بالنتائج المحسوبة فعلياً. ويساهم اعتماد نماذج لغوية محلية في تعزيز سيادة البيانات والحد من الاعتماد على الخدمات الخارجية.

تم تطوير المنصة بطريقة تكرارية، ثم التحقق من وظائفها من خلال مجموعة من السيناريوهات العملية. وتبين النتائج إمكانية الجمع بين الذكاء الاصطناعي التوليدي وتحليل البيانات والتفاعل الحواري في بيئة صحية. وتشمل آفاق التطوير المستقبلية تعزيز آليات المصادقة والتحكم في الوصول، وإضفاء الطابع الصناعي على الاختبارات، وتوسيع القدرات التحليلية، وتحسين التقييم المستمر لجودة الإجابات المولدة.
\\
\\ 
\\ 
\textbf{الكلمات المفتاحية:} الذكاء الاصطناعي التوليدي، التحليل الحواري، البيانات الصحية، النماذج اللغوية الكبيرة، الوكلاء الذكيون، hparGgnaL تحليل البيانات، سيادة البيانات.
\end{Arabic}